\section{Introduction}
\label{sec:introduction}

Partial differential equations govern phenomena across fluid dynamics, quantum
mechanics, and continuum mechanics, yet obtaining analytical solutions remains
one of the most difficult tasks in applied mathematics.  When closed-form
solutions do exist, they provide qualitative insight, asymptotic guarantees,
and computational efficiency that no numerical approximation can match.
The recent success of deep-learning-based PDE solvers has dramatically
expanded the class of problems that can be tackled computationally, but it has
done so at the cost of interpretability: a trained physics-informed neural
network (PINN)~\cite{raissi2019} produces a set of floating-point weights from
which no human-readable formula can be extracted.

\paragraph{Neural PDE solvers: accuracy without insight.}
PINNs and their descendants embed the governing PDE into the loss function of a
neural network and optimize over the network parameters so that the learned
map approximately satisfies the differential equation together with its
boundary and initial conditions.  This paradigm is flexible---it requires
neither a mesh nor a specialized numerical scheme---and has yielded
impressive results on forward and inverse problems alike.  However, the output
of such a solver is an opaque parameterization: one obtains a weight matrix,
not a formula.  Post-hoc interpretability techniques can highlight which input
regions matter, but they cannot recover the symbolic structure that a
mathematician would recognize as a solution.

\paragraph{Symbolic regression: powerful but limited in scope.}
An alternative line of work seeks to discover symbolic expressions directly from
data.  Tools such as PySR~\cite{cranmer2023} and AI~Feynman~\cite{udrescu2020}
have achieved striking success on algebraic physical laws---rediscovering
conservation laws, constitutive relations, and dimensionless
groupings from noisy measurements.  These methods search over a grammar of
elementary operations and fit scalar expressions of moderate complexity.
Yet PDE solutions pose a qualitatively different challenge.  A typical
analytical solution is not a single algebraic expression in one variable but a
\emph{product of functions of different variables}---for instance,
$\sin(x)\cos(y)\exp(-\nu t)$---and the combinatorial explosion of
multi-variable compositions quickly overwhelms generic symbolic regression
engines.

\paragraph{The EML function as a universal basis.}
Odrzywo{\l}ek~\cite{odrzywołek2026} recently proved that the two-argument function
\begin{equation}
  \label{eq:eml}
  \mathrm{eml}(x,y) \;=\; e^{x} - \ln y
\end{equation}
is \emph{universal}: every computable function can be represented as a finite
composition of $\mathrm{eml}$ with itself.  This result is striking because it
reduces the space of all computable functions to a single primitive operation,
in contrast to the multi-operation libraries assumed by conventional symbolic
regression.  In principle, one could search over EML expression trees to
discover any target function.  In practice, however, the raw EML representation
is unwieldy.  Multiplication alone---$x \cdot y$---requires a reverse-Polish
representation of length~41 when expanded purely in terms of $\mathrm{eml}$.
Directly optimizing over such trees is computationally intractable for
expressions of the complexity encountered in PDE solutions.

\paragraph{Our approach: verified primitives and compositional search.}
The key observation underlying our method is that one need not optimize at the
level of raw EML trees.  Instead, we first derive \emph{verified constructions}
of the standard elementary functions---$\exp$, $\ln$, $\sin$, $\cos$---from
the base EML operation using algebraic identities.  Each construction is
validated symbolically so that it reproduces the target function exactly.
These verified primitives then serve as differentiable building blocks in a
higher-level architecture we call the \textsc{ComposeHead}.

The \textsc{ComposeHead} represents candidate solutions as trainable linear
combinations of \emph{separable terms}, where each term is a product of
axis-aligned primitive applications:
\begin{equation}
  \label{eq:composehead}
  \hat{u}(\mathbf{x}) \;=\; \sum_{k=1}^{K} c_k \prod_{d=1}^{D}
    f_{k,d}\!\bigl(\alpha_{k,d}\, x_d + \beta_{k,d}\bigr),
\end{equation}
with each factor $f_{k,d}$ drawn from the primitive library
$\{\sin,\cos,\exp,\ln,\mathrm{id}\}$ and each $\alpha_{k,d}$, $\beta_{k,d}$,
$c_k$ a trainable scalar coefficient.  This separable structure directly
mirrors the form of most known analytical PDE solutions, in which each spatial
or temporal variable appears inside its own elementary function.  The
optimization problem thus reduces from ``find 60 arbitrary neural-network
weights'' to ``identify which functions compose, with what coefficients.''

\paragraph{From learned parameters to exact symbolic expressions.}
After training, the learned coefficients are real-valued and carry the
accumulated artifacts of gradient-based optimization---arbitrary phase offsets in
trigonometric terms, exponential biases absorbed into multiplicative constants,
and small numerical residuals.  We introduce a \emph{normalization and snapping}
pipeline that (i)~absorbs exponential biases into the linear coefficients~$c_k$,
(ii)~canonicalizes trigonometric phases by exploiting the identities
$\sin(x + \pi/2) = \cos(x)$ and $\cos(x + \pi) = -\cos(x)$, and
(iii)~rounds the resulting coefficients to nearby rational or otherwise
recognizable values.  The output is a clean symbolic expression that can be
verified independently against the governing PDE.

\paragraph{Scope and limitations.}
We emphasize that the method proposed here is best understood as
\emph{symbolic regression specialized to PDE solutions}, not as a general PDE
solver.  The separable ansatz in Equation~\eqref{eq:composehead} is expressive
for the broad class of problems whose analytical solutions factor across
coordinate dimensions, but it does not address existence, uniqueness, or
regularity questions, nor does it replace mesh-based solvers for problems
lacking closed-form solutions.  Our claims are accordingly conservative: we
demonstrate that, when an analytical solution exists and has separable structure,
the method can recover it exactly from PDE-residual supervision alone.

\paragraph{Contributions.}
The principal contributions of this work are as follows.
\begin{itemize}
  \item \textbf{Verified EML primitive constructions.}  We provide
    differentiable, algebraically verified implementations of $\exp$, $\ln$,
    $\sin$, and $\cos$ built entirely from compositions of the EML function,
    establishing a concrete link between the universality theorem
    of~\cite{odrzywołek2026} and practical symbolic computation.

  \item \textbf{ComposeHead architecture.}  We introduce a neural architecture
    whose hypothesis class is the set of linear combinations of separable
    products of axis-aligned elementary functions, enabling gradient-based
    search over symbolic PDE solution candidates.

  \item \textbf{Normalization and snapping pipeline.}  We describe a
    post-training procedure that absorbs phase offsets and exponential biases
    into canonical form, then snaps coefficients to exact values, converting
    a continuously parameterized model into a closed-form symbolic expression.

  \item \textbf{Exact recovery of the Taylor--Green vortex.}  We demonstrate
    the full pipeline on the two-dimensional incompressible Navier--Stokes
    equations, recovering the velocity field
    $u = \sin(x)\cos(y)\exp(-t/5)$,\;$v = -\cos(x)\sin(y)\exp(-t/5)$
    with pointwise error below $10^{-6}$, confirming that the discovered
    symbolic expression is an exact solution of the governing equations.
\end{itemize}

\noindent The remainder of the paper is organized as follows.
Section~\ref{sec:background} reviews the EML universality result and the
relevant symbolic regression literature.
Section~\ref{sec:method} details the primitive constructions, the
\textsc{ComposeHead} architecture, and the normalization pipeline.
Section~\ref{sec:experiments} presents the Taylor--Green vortex experiment and
ablation studies.
Section~\ref{sec:discussion} discusses limitations and directions for future
work.
